\chapter[1959 --- Ochoa et al.]{1959: Severo Ochoa, Arthur Kornberg}
\label{ch:1959_Ochoa}

\section*{Main Contribution}

\textit{Discovered the mechanisms by which cells synthesize DNA and RNA, explaining how genetic information is copied and expressed.}

\section{Official Citation}

for their discovery of the mechanisms in the biological synthesis of ribonucleic acid and deoxyribonucleic acid

\section{Background and Historical Context}

The 1959 Nobel Prize in Physiology or Medicine was awarded to Severo Ochoa and Arthur Kornberg for their discovery of the mechanisms in the biological synthesis of ribonucleic acid (RNA) and deoxyribonucleic acid (DNA). This prize recognized two related breakthroughs that provided the first direct evidence for the enzymatic synthesis of nucleic acids---the molecules that carry genetic information in all living organisms.

\subsection{The Central Question of Molecular Biology}

By the late 1950s, the structure of DNA had been elucidated by James Watson and Francis Crick, and researchers had begun to understand the role of DNA as the carrier of genetic information. However, the mechanisms by which DNA was replicated and by which genetic information was transferred from DNA to RNA to proteins were still unknown. The work of Ochoa and Kornberg provided the first tools for studying these fundamental processes.

\subsection{Severo Ochoa}

Severo Ochoa was born on September 24, 1905, in Luarca, Spain. He studied medicine at the University of Madrid, where he developed an interest in biochemistry and physiology. After receiving his medical degree in 1929, Ochoa worked at the University of Madrid and later at the Kaiser Wilhelm Institute for Medical Research in Heidelberg, Germany, where he studied the metabolism of carbohydrates and nucleic acids.

In 1938, Ochoa emigrated to the United States to escape the Spanish Civil War. He joined the faculty at Washington University in St. Louis and later at New York University, where he would spend most of his career. Ochoa was a skilled biochemist who was interested in the mechanisms of energy metabolism and nucleic acid synthesis.

\subsection{Arthur Kornberg}

Arthur Kornberg was born on March 3, 1918, in Brooklyn, New York. He studied at the College of the City of New York, where he earned his bachelor's degree in 1937, and at the University of Rochester School of Medicine, where he earned his medical degree in 1941. After completing his medical training, Kornberg worked at the National Institutes of Health (NIH), where he developed his interest in the enzymatic synthesis of nucleic acids.

Kornberg was a meticulous and persistent researcher who was determined to isolate the enzymes responsible for DNA synthesis. His work on DNA polymerase would become one of the major discoveries in the history of molecular biology.

\subsection{The Nucleic Acid Synthesis Puzzle}

By the late 1950s, researchers had established that DNA was the carrier of genetic information and that RNA played an important role in protein synthesis. However, the mechanisms by which DNA and RNA were synthesized were still unknown. The work of Ochoa and Kornberg would provide the first direct evidence for the enzymatic synthesis of these molecules and would lay the foundation for the development of modern molecular biology.

\section{The Discovery/Contribution}

The work of Ochoa and Kornberg on nucleic acid synthesis involved the discovery and characterization of enzymes that could synthesize RNA and DNA from their building blocks.

\subsection{Ochoa's Work on Polynucleotide Phosphorylase}

Severo Ochoa's primary contribution was the discovery of polynucleotide phosphorylase, an enzyme that could catalyze the synthesis of RNA from nucleotide precursors. Ochoa's discovery came about through his study of the metabolism of nucleotides in bacteria.

\subsubsection{The Discovery}

In 1955, Ochoa and his colleagues discovered that extracts of the bacterium \textit{Azotobacter vinelandii} could catalyze the synthesis of RNA from nucleotide diphosphates (NDPs) in the presence of inorganic phosphate. The enzyme responsible for this reaction was named polynucleotide phosphorylase (PNPase).

Ochoa showed that PNPase could catalyze the following reaction:
\[ \text{NDP} \rightleftharpoons \text{RNA} + \text{Pi} \]

This reaction is reversible, and the direction of the reaction depends on the concentrations of the reactants and products.

\subsubsection{Significance of the Discovery}

The discovery of PNPase was significant for several reasons:

\begin{itemize}
\item It provided the first evidence that RNA could be synthesized enzymatically
\item It demonstrated that enzymes could catalyze the formation of phosphodiester bonds between nucleotides
\item It provided a tool for synthesizing RNA in the laboratory, which was essential for studies of the genetic code
\end{itemize}

However, it was later discovered that PNPase is not the enzyme responsible for RNA synthesis in living cells. Instead, it plays a role in RNA degradation and processing. The enzyme responsible for RNA synthesis in vivo is RNA polymerase, which was discovered later.

\subsection{Kornberg's Work on DNA Polymerase}

Arthur Kornberg's primary contribution was the discovery of DNA polymerase, the enzyme responsible for DNA synthesis. Kornberg's discovery was the result of a systematic search for the enzyme or enzymes responsible for DNA replication.

\subsubsection{The Experimental Approach}

Kornberg's approach was to purify the enzymes responsible for DNA synthesis from extracts of \textit{Escherichia coli}. He developed a sensitive assay for DNA synthesis that measured the incorporation of radioactive nucleotides into DNA. Using this assay, he was able to follow the activity of the enzyme through successive steps of purification.

\subsubsection{The Discovery of DNA Polymerase}

In 1956, Kornberg and his colleagues isolated an enzyme they called DNA polymerase (later renamed DNA polymerase I) from E. coli. This enzyme could catalyze the synthesis of DNA from deoxynucleoside triphosphates (dNTPs) in the presence of a DNA template. The enzyme required all four dNTPs (dATP, dGTP, dCTP, and dTTP) as well as a DNA template to synthesize DNA.

Kornberg showed that DNA polymerase could replicate DNA in vitro (in a test tube) and that the newly synthesized DNA was complementary to the template DNA. This was a landmark discovery because it provided the first direct evidence for the enzymatic synthesis of DNA.

\subsubsection{Characteristics of DNA Polymerase}

Kornberg's work on DNA polymerase revealed several important properties of the enzyme:

\begin{itemize}
\item The enzyme requires a DNA template to direct the synthesis of new DNA
\item The enzyme adds nucleotides to the 3' end of the growing DNA chain
\item The enzyme requires all four dNTPs as substrates
\item The enzyme has a proofreading function that allows it to correct errors in the newly synthesized DNA
\end{itemize}

\subsubsection{Later Discoveries}

Kornberg's work on DNA polymerase led to the discovery of additional DNA polymerases in E. coli and other organisms. It was eventually discovered that DNA polymerase I is not the primary enzyme responsible for DNA replication in E. coli. Instead, it plays a role in DNA repair and in the removal of RNA primers during DNA replication. The primary enzyme responsible for DNA replication in E. coli is DNA polymerase III, which was discovered later.

\section{Impact on Modern Medicine}

The work of Ochoa and Kornberg on nucleic acid synthesis has had a significant impact on modern medicine. Their discoveries have contributed to our understanding of genetic diseases, cancer, and infectious diseases, and have led to the development of numerous medical technologies and therapies.

\subsection{Understanding DNA Replication}

Kornberg's discovery of DNA polymerase provided the first direct evidence for the enzymatic synthesis of DNA and laid the foundation for our understanding of DNA replication. The mechanisms of DNA replication that were elucidated through Kornberg's work are now used to understand a wide range of genetic diseases and to develop new therapies.

\subsection{Polymerase Chain Reaction (PCR)}

Kornberg's work on DNA polymerase also contributed to the development of the polymerase chain reaction (PCR), a technique that allows researchers to amplify specific DNA sequences. PCR, developed by Kary Mullis in the 1980s, has revolutionized molecular biology and has numerous applications in medicine, including:

\begin{itemize}
\item Diagnosis of genetic diseases
\item Detection of infectious diseases
\item Forensic identification
\item Prenatal diagnosis of genetic disorders
\end{itemize}

\subsection{Understanding RNA Synthesis}

Ochoa's work on polynucleotide phosphorylase contributed to our understanding of RNA metabolism and processing. Although PNPase is not the enzyme responsible for RNA synthesis in vivo, Ochoa's work laid the foundation for the discovery of RNA polymerase and other enzymes involved in RNA metabolism.

\subsection{Nucleic Acid-Based Therapies}

The understanding of nucleic acid synthesis that Ochoa and Kornberg helped to establish has also contributed to the development of nucleic acid-based therapies. These include:

\begin{itemize}
\item Antisense oligonucleotides: Synthetic DNA or RNA molecules that can bind to specific mRNA molecules and block their translation into proteins
\item RNA interference (RNAi): A technique that uses small RNA molecules to silence specific genes
\item Gene therapy: A technique that involves introducing functional genes into cells to correct genetic defects
\end{itemize}

\subsection{Diagnostic Applications}

The understanding of nucleic acid synthesis has also contributed to the development of diagnostic tests for genetic diseases and infectious diseases. These tests can detect specific DNA or RNA sequences in patient samples and are used to diagnose a wide range of conditions.

\section{Legacy}

The legacy of Severo Ochoa and Arthur Kornberg extends far beyond their Nobel Prize-winning work on nucleic acid synthesis. Their discoveries laid the foundation for modern molecular biology and have contributed to numerous advances in medicine and biotechnology.

\subsection{Foundation for Molecular Biology}

The work of Ochoa and Kornberg provided the first tools for studying DNA and RNA synthesis in the laboratory. These tools were essential for the subsequent elucidation of the genetic code, the development of recombinant DNA technology, and the sequencing of the human genome. The principles established by Ochoa and Kornberg continue to guide research in molecular biology today.

\subsection{Advances in Medicine}

The understanding of nucleic acid synthesis that Ochoa and Kornberg helped to establish has contributed to numerous advances in medicine, including the development of diagnostic tests for genetic diseases, the understanding of cancer biology, and the development of nucleic acid-based therapies. These advances have improved the lives of millions of patients and continue to drive progress in medicine.

\subsection{Biotechnology Applications}

The enzymes and techniques developed by Ochoa and Kornberg have also found numerous applications in biotechnology. DNA polymerase is now used in a wide range of biotechnological applications, including PCR, DNA sequencing, and genetic engineering. These technologies have revolutionized the biotechnology industry and have contributed to the development of new products and therapies.

\subsection{Recognition and Honors}

In addition to the Nobel Prize, Ochoa and Kornberg received numerous other honors and awards for their work. Ochoa received the National Medal of Science in 1979 and was elected to the National Academy of Sciences. Kornberg received the National Medal of Science in 1979 and was elected to the National Academy of Sciences.

Their contributions to molecular biology and medicine are remembered and celebrated by researchers and clinicians around the world. Their discoveries continue to influence our understanding of nucleic acid synthesis and have led to numerous advances in medicine and biotechnology, ensuring that their legacy will endure for generations to come.


\section{Modern Developments}

Ochoa and Kornberg's work on nucleic acid synthesis has enabled modern molecular biology. DNA polymerases and RNA polymerases are essential tools for PCR, gene expression studies, and vaccine production. Understanding of DNA replication has led to antiviral drugs (acyclovir for herpes, tenofovir for HIV) that target viral DNA synthesis. PCR technology, which revolutionized diagnostics and forensics, directly builds on their enzymatic discoveries.


\section{Bibliography}

\begin{thebibliography}{9}
\bibitem{nobel-1959}
Nobel Prize Outreach, ``The Nobel Prize in Physiology or Medicine 1959,''\emph{NobelPrize.org}.\\
\url{https://www.nobelprize.org/prizes/medicine/1959/summary/}

\bibitem{lecture-1959}
Severo Ochoa, ``Nobel Lecture,''\emph{NobelPrize.org}. Winner-authored primary source; downloadable from the official Nobel Prize archive.\\
\url{https://www.nobelprize.org/prizes/medicine/1959/ochoa/lecture/}
\end{thebibliography}
